\documentclass{diploma}

% Текст-заполнитель
\usepackage{lipsum}
\SetLipsumText{fishes} % fishes.ltd.tex

\student{Гончаров Андрей Яковлевич}
\group{М8О-407Б-19}
\theme{Создание отчёта по ВКР}

\supervisor{Орлов Марсель Максимович}
\firstConsultant{Алешина Анна Константиновна}
\secondConsultant{Свиридов Михаил Артемьевич}
\reviewer{Ильина Алина Даниловна}

\faculty{№ 8 <<Компьютерные науки и прикладная математика>>}
\department{806}
\speciality{01.02.03 <<Прикладная математика и информатика>>}
\profile{Информатика}

% Здесь необходим разрыв строки из-за особенностей титульной страницы
\departmentFullName{№ 806 <<Вычислительная математика \\ и программирование>>}
\headOfDepartment{Крылов Сергей Сергеевич}

% Дата. Оставляем пустое место для дня
\date{\uline{\hspace{24pt}} мая \the\year\ года}

\newacronym{cd}{CD}{compact disk}

\newglossaryentry{id1}{ % Нужны разные id, можно ставить просто последовательно
    name={\LaTeX},
    description={система компьютерной вёрстки} 
}
\newglossaryentry{id2}{
    name={МАИ},
    description={Московский Авиационный Институт} 
}
\newglossaryentry{id3}{
    name={$S^{2}$},
    description={формула в списке обозначений} 
}


\addbibresource{main.bib}

% Иллюстрации всегда по центру
\makeatletter
\g@addto@macro\@floatboxreset\centering
\makeatother
    
\begin{document}
    \maketitle

    % \includepdf[pages=-]{extra/task} % Задание
    \setcounter{page}{2} % Устанавливает счётчик страниц

    \abstract % Структурный элемент: РЕФЕРАТ

КЛЮЧЕВЫЕ СЛОВА, КЛЮЧЕВЫЕ СЛОВА, КЛЮЧЕВЫЕ СЛОВА, КЛЮЧЕВЫЕ СЛОВА, КЛЮЧЕВЫЕ СЛОВА

\lipsum[1-2]
 % Реферат

    \tableofcontents % Содержание 
    \termsanddefenitions % Термины и определения
    \listofabbreviations % Перечень сокращений и обозначений
 
    \introduction % Структурный элемент: ВВЕДЕНИЕ

\lipsum[1-5]

Цитирование источника 1 \cite{Wikipedia1}.
 % Введение

    % Название разделов -- все прописные
\section{КОМПИЛЯЦИЯ}

Шаблон предназначен для сборки с использованием \texttt{latexmk}. Рекомендуемый компилятор --- \texttt{xelatex}:

\begin{verbatim}
latexmk -xelatex main.tex
\end{verbatim}

Директория \texttt{diploma} должна быть достижима от основной директории отчета (скопирована локально или, например, как \texttt{git submodule}). В файле \texttt{.latexmkrc} в переменной \texttt{TEXINPUTS} задается относительный путь от директории отчета до директории \texttt{diploma} (\texttt{//} в конце пути обязательны).

\begin{verbatim}
ensure_path( 'TEXINPUTS', '..//');
\end{verbatim}

Или можете использовать шаблон в overleaf \url{https://www.overleaf.com/read/vyqpdcfnhmsy}.

\lipsum[1-2] % Основная часть
    \section{Использование титульной страницы}

Заполнение полей показано в файле \texttt{main.tex}, приведу ещё раз здесь. 
Желательно заполнить все поля по образцу, чтобы не было проблем с тем, 
что что-то может быть не определено.

\begingroup
\fontsize{10pt}{12pt}\selectfont
\begin{verbatim}
\student{Гончаров Андрей Яковлевич}
\group{М8О-407Б-19}
\theme{Создание отчёта по ВКР}

\supervisor{Орлов Марсель Максимович}
\firstConsultant{Алешина Анна Константиновна}
\secondConsultant{Свиридов Михаил Артемьевич}
\reviewer{Ильина Алина Даниловна}

\faculty{№ 8 <<Компьютерные науки и прикладная математика>>}
\department{806}
\speciality{01.02.03 <<Прикладная математика и информатика>>}
\profile{Информатика и компьютерные науки}

% Здесь необходим разрыв строки из-за особенностей титульной страницы
\departmentFullName{№ 806 <<Вычислительная математика \\ и программирование>>}
\headOfDepartment{Крылов Сергей Сергеевич}
\date{XX мая}
\end{verbatim}
\endgroup

\lipsum[3]

Цитирование источника 2 \cite{Wikipedia2}.
    \section{ТЕХНИЧЕСКОЕ ЗАДАНИЕ И КАЛЕНДАРНЫЙ ПЛАН}

Вы не поверите, но это вообще отдельные сущности. Они не должны быть
в отчете. Требуется именно такая сдачи --- с КП и ИЗ, вложенными в РПЗ.

\lipsum[1-2]

Цитирование источника 3 \cite{Wikipedia3}.

    \section{Создание реферата}

Первая строка там уже есть --- с использованием счётчиков страниц, рисунков, 
таблиц, источников и приложений. 
Необходимо заполнить ключевые слова и краткий обзор своей работы, 
куда входит цель, этапы и прочее...

\lipsum[3-4]

Цитирование источника 4 \cite{Wikipedia4} \cite{cite_1_2} \cite{cite_1_15} \cite{cite_1_16}.

    \section{ЗАДАНИЕ СПИСКА ТЕРМИНОВ, СОКРАЩЕНИЙ И ОПРЕДЕЛЕНИЙ}

Определения необходимо записывать в файл \texttt{terms.tex}, сокращения --- в \texttt{glossary.tex}.
Файлы необходимо подключать до начала документа, при помощи \texttt{\textbackslash input}.

Обратите внимание, что термины и сокращения необходимо приводить в алфавитном порядке, пока что это можно сделать только вручную.

Для того чтобы списки подключились в качестве части отчёта, 
нужно использовать команды \texttt{\textbackslash termsanddefenitions} и \texttt{\textbackslash listofabbreviations} 
сразу после команды \texttt{\textbackslash tableofcontents}.

    \section{Использование рисунков}

Вставляются рисунки как обычно --- через 
\texttt{\textbackslash includegraphics} и окружение 
\texttt{figure}. 
Пожалуйста, не используйте \texttt{[H]} --- 
в этом шаблоне уже настроена среда картинок так, 
что она вставится как можно ближе к тексту по возможности. 
Использование \texttt{[H]} приводит к большим и некрасивым 
разрывам текста. Это же касается и таблиц.

На рисунке~\ref{fig:fig01} показан логотип МАИ. 
Также тут видно, что ссылка на рисунок работает.
\begin{figure}
  \centering
  \includegraphics[scale=0.6]{inc/mai}
  \caption{Логотип МАИ}
  \label{fig:fig01}
\end{figure}


\lipsum[3]

Цитирование источника 7 \cite{Wikipedia7}.

    \section{Цитирование источников}

Вот так \cite{Article} можно цитировать статьи. 
Заполнение представлено в файле \texttt{main.bib}. 
Пожалуйста, указывайте \texttt{russian} в качестве параметра \texttt{language}!

Аналогично можно цитировать сайты в интернете, но нужно будет добавить 
дату обращения \cite{Wikipedia}.

Также можно вставлять ссылки командой \url{https://youtu.be/dQw4w9WgXcQ}, например.

Больше примеров для оформления библиографических ссылок можно найти в 
документации к пакету \texttt{gost2008}: 
\url{http://tug.ctan.org/tex-archive/biblio/bibtex/contrib/gost/doc/examples/cp1251/gost2008.pdf}.

\lipsum[1-2]

Цитирование источника 2 \cite{cite_1_11}.

    \section{Работа с таблицами}

Далее рассматриваются варианты создания таблиц. 
На этом моменте лучше смотреть в исходник.

\lipsum[1]

\subsection{Использование tabular}

Самый простой и стандартный способ --- использование tabular. 
Так мы создали таблицу~\ref{tab:tab1}. 
Обратите внимание на использование \texttt{hhline} --- этот пакет 
позволяет подчёркивать не всю линию, а только те столбцы, которые 
нужно подчеркнуть. Полезно при использовании мультистрок или мультистолбцов.

\begin{table}    
    \caption{Пример короткой таблицы с использованием tabular}
    \begin{tabular}{|r|c|c|c|l|}\hline
    Тело      & $F$ & $V$  & $E$ & $F+V-E-2$ \\ \hline
    Тетраэдр  & 4   & 4    & 6   & 0         \\ \hhline{~-~-~}
    Куб       & 6   & 8    & 12  & 0         \\ \hhline{--~~~}
    Октаэдр   & 8   & 6    & 12  & 0         \\ \hhline{-----}
    Додекаэдр & 20  & 12   & 30  & 0         \\ \hline
    Икосаэдр  & 12  & 20   & 30  & 0         \\ \hline
    \end{tabular}
    \label{tab:tab1}
\end{table}

\lipsum[1]

\subsection{Использование tabularx}

Этот способ мне нравится больше, потому что позволяет задать ширину 
таблицы~\ref{tab:tab2}. Единственное ограничение --- в этой таблице должен 
быть столбец, обозначаемый \texttt{X}, который может быть растянут до 
нужного размера (заполнения листа).

\begin{table}    
    \caption{Пример короткой таблицы с tabularx}
    \begin{tabularx}{\textwidth}{|X|c|c|c|l|}\hline
    Тело      & $F$ & $V$  & $E$ & $F+V-E-2$ \\ \hline
    Тетраэдр  & 4   & 4    & 6   & 0         \\ \hhline{~-~-~}
    Куб       & 6   & 8    & 12  & 0         \\ \hhline{--~~~}
    Октаэдр   & 8   & 6    & 12  & 0         \\ \hhline{-----}
    Додекаэдр & 20  & 12   & 30  & 0         \\ \hline
    Икосаэдр  & 12  & 20   & 30  & 0         \\ \hline
    \end{tabularx}
    \label{tab:tab2}
\end{table}

\lipsum[1]

\subsection{Использование longtable}

\texttt{longtable} используется для задания многостраничных длинных таблиц. 
Необходимо указать надпись, которая будет над таблицей на следующем листе. 
Пример ниже, в таблице~\ref{tab:tab3}. Обратите внимание --- здесь уже не 
нужно использовать среду \texttt{table}.

 
\begin{longtable}[H]{|c|c|c|c|c|}
    \caption{Пример использования длинной таблицы на несколько листов, 
    а также пример использования длинного заголовка таблицы}\label{tab:tab3}\\ \hline
    \endfirsthead
    \caption*{Продолжение таблицы \ref{tab:tab3}}\\ \hline
    \endhead
    \hline
    \endfoot
    \hline
    \endlastfoot
    
    Тело      & $F$ & $V$  & $E$ & $F+V-E-2$ \\ \hline
    Тетраэдр  & 4   & 4    & 6   & 0         \\ \hhline{~-~-~}
    Куб       & 6   & 8    & 12  & 0         \\ \hhline{--~~~}
    Октаэдр   & 8   & 6    & 12  & 0         \\ \hhline{-----}
    Додекаэдр & 20  & 12   & 30  & 0         \\ \hline
    Икосаэдр  & 12  & 20   & 30  & 0         \\ \hline
    Тело      & $F$ & $V$  & $E$ & $F+V-E-2$ \\ \hline
    Тетраэдр  & 4   & 4    & 6   & 0         \\ \hhline{~-~-~}
    Куб       & 6   & 8    & 12  & 0         \\ \hhline{--~~~}
    Октаэдр   & 8   & 6    & 12  & 0         \\ \hhline{-----}
    Додекаэдр & 20  & 12   & 30  & 0         \\ \hline
    Икосаэдр  & 12  & 20   & 30  & 0         \\ \hline
    Тело      & $F$ & $V$  & $E$ & $F+V-E-2$ \\ \hline
    Тетраэдр  & 4   & 4    & 6   & 0         \\ \hhline{~-~-~}
    Куб       & 6   & 8    & 12  & 0         \\ \hhline{--~~~}
    Октаэдр   & 8   & 6    & 12  & 0         \\ \hhline{-----}
    Додекаэдр & 20  & 12   & 30  & 0         \\ \hline
    Икосаэдр  & 12  & 20   & 30  & 0         \\ \hline
    Тело      & $F$ & $V$  & $E$ & $F+V-E-2$ \\ \hline
    Тетраэдр  & 4   & 4    & 6   & 0         \\ \hhline{~-~-~}
    Куб       & 6   & 8    & 12  & 0         \\ \hhline{--~~~}
    Октаэдр   & 8   & 6    & 12  & 0         \\ \hhline{-----}
    Додекаэдр & 20  & 12   & 30  & 0         \\ \hline
    Икосаэдр  & 12  & 20   & 30  & 0         \\ \hline
    Тело      & $F$ & $V$  & $E$ & $F+V-E-2$ \\ \hline
    Тетраэдр  & 4   & 4    & 6   & 0         \\ \hhline{~-~-~}
    Куб       & 6   & 8    & 12  & 0         \\ \hhline{--~~~}
    Октаэдр   & 8   & 6    & 12  & 0         \\ \hhline{-----}
    Додекаэдр & 20  & 12   & 30  & 0         \\ \hline
    Икосаэдр  & 12  & 20   & 30  & 0         \\ \hline
    Тело      & $F$ & $V$  & $E$ & $F+V-E-2$ \\ \hline
    Тетраэдр  & 4   & 4    & 6   & 0         \\ \hhline{~-~-~}
    Куб       & 6   & 8    & 12  & 0         \\ \hhline{--~~~}
    Октаэдр   & 8   & 6    & 12  & 0         \\ \hhline{-----}
    Додекаэдр & 20  & 12   & 30  & 0         \\ \hline
    Икосаэдр  & 12  & 20   & 30  & 0         \\ \hline
    Тело      & $F$ & $V$  & $E$ & $F+V-E-2$ \\ \hline
    Тетраэдр  & 4   & 4    & 6   & 0         \\ \hhline{~-~-~}
    Куб       & 6   & 8    & 12  & 0         \\ \hhline{--~~~}
    Октаэдр   & 8   & 6    & 12  & 0         \\ \hhline{-----}
    Додекаэдр & 20  & 12   & 30  & 0         \\ \hline
    Икосаэдр  & 12  & 20   & 30  & 0         \\ \hline
\end{longtable}


\lipsum[1]

Цитирование источника 2 \cite{Article3}.

    \section{ПЕРЕЧИСЛЕНИЯ}

По ГОСТ перечисления начинаются с букв.

Добавим немного текста. \lipsum[1][1]

\begin{itemize}
    \item первый элемент перечисления,
    \item второй элемент перечисления,
    \item третий элемент перечисления.
\end{itemize}

\lipsum[1][1]

\begin{enumerate}
    \item перечисление с номерами;
    \item номера первого уровня (да, ГОСТ требует именно так --- сначала буквы, 
    на втором уровне --- цифры)
    (да, на первом уровне выравнивание элементов как у обычных абзацев);
        \begin{enumerate}
            \item номера второго уровня;
            \item номера второго уровня;
        \end{enumerate}
    \item проверяем то;
    \item что нужные буквы;
    \item отсутствуют;
    \item в нумерации;
    \item ё нет;
    \item з тоже нет;
    \item последний элемент списка.
\end{enumerate}

Цитирование источника 2 \cite{Article3}.

    \section{ФОРМУЛЫ}

Можно сделать заинлайненую формулу: $E = mc^2$. Можно сделать формулу по центру (окружение \texttt{equation}):

\begin{equation}
    \label{f1}
    E = mc^2,
\end{equation}

где $E$ --- энергия,\par $m$ --- масса,\par $c$ --- скорость света.

В таком случае точки и знаки препинания лучше оставить внутри формулы. Обратите внимание, что пояснения указываются в порядке появления в формуле, каждое пояснение с новой строки (для этого используйте \texttt{\textbackslash par} после каждого пояснения). \textbf{Необходимо} оставлять пустую строку до (и после) окружений \texttt{equation}, \texttt{align}, \texttt{gather}, \texttt{split}, etc. иначе не будет нужного пустого расстояния.  

Также можно ссылаться на формулу \ref{f1}.

Несколько формул в окружении \texttt{align}:

\begin{align}
x'_1 &= \frac{-3 x_1 + 4 x_2 - x_3}{2 \tau} \\
x'_i &= \frac{x_{i+1} - x_{i-1}}{2 \tau}, \qquad i = \overline{2, n-1} \\
x'_n &= \frac{x_{n-2} - 4 x_{n-1} + 3 x_n}{2 \tau}
\end{align}

Несколько формул в окружении \texttt{gather}:

\begin{gather} 
2x - 5y =  8 \\ 
3x^2 + 9y =  3a + c
\end{gather}

Формула на нескольких строках в окружении \texttt{split}. По умолчанию здесь тоже будет большое расстояние между строками. Чтобы этого избежать используйте команду \texttt{\textbackslash setlength\{\textbackslash jot\}\{3pt\}} внутри \texttt{equation}:

\begin{equation}
\setlength{\jot}{3pt}
\begin{split}
F = {} &\{F_{x} \in  F_{c} : (|S| > |C|) \cap {} \\
& \cap (\mathrm{minPixels}  < |S| < \mathrm{maxPixels}) \cap {} \\
& \cap (|S_{\mathrm{conected}}| > |S| - \varepsilon) \}
\end{split}
\end{equation}
    \section{МАТЕМАТИКА: ТЕОРЕМЫ, ПРИМЕРЫ, ОПРЕДЕЛЕНИЯ И ЛЕММЫ}

\begin{definition}\label{dfn1}
    Это определение и оно нумеруется сквозной нумерацией по всему документу.
\end{definition}

\begin{theorem}\label{th1}
    Это теорема и она также имеет сквозную нумерацию. Ссылка на определение~\ref{dfn1}.
\end{theorem}

\begin{proof}
    Это доказательство.
\end{proof}

\begin{corollary}\label{cor1}
    Следствие имеет нумерацию в пределах одной теоремы. Ссылка на теорему~\ref{th1}.
\end{corollary}

\begin{example}\label{ex1}
    Пример также можно приводить в стиле теоремы. Нумерация сквозная. Ссылка на следствие~\ref{cor1}.
\end{example}

Ссылка на пример \ref{ex1}.
    \section{ОПИСАНИЕ АЛГОРИТМОВ}

Псевдокод необходимо вставлять как иллюстрации --- он должен быть помещен в окружение \texttt{figure}, подписан и иметь ссылку в тексте.

Также из-за особенностей \texttt{algorithm}, для него обязательна опция \texttt{H} и внутри \texttt{figure} его еще надо оборачивать в \texttt{minipage} размером с длину строки

Пример использования алгоритма представлен на рисунке~\ref{alg:alg1}:

\begin{figure}
\begin{minipage}{\linewidth}
\begin{algorithm}[H]  
    \SetAlgoVlined
    \KwData{Входные данные}
    \KwResult{Как прочитать книгу }
    инициализация\;
    \While{есть непонятая глава книги}{
        прочитать текущую главу\;
        \eIf{понятно}{
            перейти к следующей главе\;
            текущей главой становится следующая глава\;
        }{
            перейти к началу текущей главы\;
        }
    } 
\end{algorithm}
\end{minipage}
\caption{Как прочитать книгу}
\label{alg:alg1}
\end{figure}

\pagebreak

Ещё один пример на рисунке~\ref{alg:generalGP}:

\begin{figure}
\begin{minipage}{\linewidth}
\begin{algorithm}[H]
	\SetAlgoVlined %% Это соединяет линиями логические части
	%% алгоритма типа if-then-else
	
	\KwData{ experiment.data} %% здесь можно указать исходные параметры
	\KwResult{ output, xoptimal } %% результат работы программы
	x=0\;
	\While{ $\tau_{norm} > \varepsilon_{tol}$ }{
		$s_{k-1} \leftarrow x_k - x_{k-1}$\;
		\tcc{Step lenght computation:} %% это комментарий, который будет виден.
		\eIf{$k$ is even}{
			$ \alpha_k^{ABB} = \frac{ s_{k-1}^T y_{k-1}}{y_{k-1}^T y_{k-1}}$\;
		}{ %% ELSE
		$\alpha_k^{ABB} = \frac{ s_{k-1}^T s_{k-1}}{s_{k-1}^T y_{k-1}}$\;
	} %% END eIf{$k$ is even}{
	$k \leftarrow k + 1$\;
	\For{ i = 1}{
		$x_{i+1} = P_\Omega(x_i - \alpha_k^{ABB}*g_k)$\;
	} %% END For{ i = 1}{
	\tcc{Compute the termination constant} %% это комментарий, который будет виден.
	$\tau_{norm} = abs ( ||x_{k}||_2 - ||x_{k-1}||_2)$\;
} %% END While{ $\tau_{norm} > \varepsilon_{tol}$ }{
\end{algorithm}
\end{minipage}
\caption{Псевдо-код алгоритма}
\label{alg:generalGP}
\end{figure}

Цитирование источника 2 \cite{cite_1_10}.

    \section{Вставка кода}

Хотя код и нужно вставлять в качестве приложений, покажу, как его можно использовать с помощью пакета \texttt{lstlisting}. Это оформлено не по ГОСТ, но я особо-то стандартов на листинги и нет. Если найдёте --- сообщите.

\listing[frame=rlbt,language=Python,caption={Пример использования листинга}]{inc/example.py}

Стоит учесть, что данный стиль листингов не по ГОСТ и подсвечивает синтаксис.

Следующий листинг показывает разрыв страницы и задание кода сразу в \texttt{tex}-файле:

\begin{lstlisting}[frame=rlbt,language=Python,caption={Длинный листинг}]
import numpy as np
    
def incmatrix(genl1,genl2):
    m = len(genl1)
    n = len(genl2)
    M = None #to become the incidence matrix
    VT = np.zeros((n*m,1), int)  #dummy variable
    
    #compute the bitwise xor matrix
    M1 = bitxormatrix(genl1)
    M2 = np.triu(bitxormatrix(genl2),1) 

    for i in range(m-1):
        for j in range(i+1, m):
            [r,c] = np.where(M2 == M1[i,j])
            for k in range(len(r)):
                VT[(i)*n + r[k]] = 1;
                VT[(i)*n + c[k]] = 1;
                VT[(j)*n + r[k]] = 1;
                VT[(j)*n + c[k]] = 1;
                
                if M is None:
                    M = np.copy(VT)
                else:
                    M = np.concatenate((M, VT), 1)
                
                VT = np.zeros((n*m,1), int)
    
    return M
\end{lstlisting}
    
    \conclusion

Товарищи! рамки и место обучения кадров играет важную роль в 
формировании форм развития. Таким образом дальнейшее развитие различных форм 
деятельности играет важную роль в формировании дальнейших направлений развития. 
Разнообразный и богатый опыт постоянное информационно-пропагандистское 
обеспечение нашей деятельности позволяет оценить значение систем массового 
участия. Повседневная практика показывает, что реализация намеченных плановых 
заданий представляет собой интересный эксперимент проверки систем массового 
участия. 

Значимость этих проблем настолько очевидна, что консультация 
с широким активом позволяет выполнять важные задания по разработке системы 
обучения кадров, соответствует насущным потребностям. Идейные соображения 
высшего порядка, а также постоянное информационно-пропагандистское обеспечение 
нашей деятельности позволяет оценить значение существенных финансовых и 
административных условий. Задача организации, в особенности же сложившаяся 
структура организации обеспечивает широкому кругу (специалистов) участие в 
формировании систем массового участия. Разнообразный и богатый опыт начало 
повседневной работы по формированию позиции представляет собой интересный 
эксперимент проверки форм развития. % Заключение

    \printbibliography % Список литературы

    \appendix % Приложения
    \appendixsection{Приложение с рисунком}
    

Заголовок приложения задаётся командой \texttt{\textbackslash appendixsection}. 
Рисунок~\ref{fig:a1} идёт с нумерацией приложения.

\begin{figure}
    \centering
    \includegraphics[width=5cm]{inc/mai}
    \caption{Рисунок в приложении}
    \label{fig:a1}
\end{figure}

\lipsum[1-2]

    \appendixsection{Приложение с таблицей}

Таблица~\ref{tab:b1} идёт с нумерацией приложения.

\begin{table}    
    \caption{Пример короткой таблицы с использованием tabular}
    \begin{tabular}{|r|c|c|c|l|}\hline
    Тело      & $F$ & $V$  & $E$ & $F+V-E-2$ \\ \hline
    Тетраэдр  & 4   & 4    & 6   & 0         \\ \hhline{~-~-~}
    Куб       & 6   & 8    & 12  & 0         \\ \hhline{--~~~}
    Октаэдр   & 8   & 6    & 12  & 0         \\ \hhline{-----}
    Додекаэдр & 20  & 12   & 30  & 0         \\ \hline
    Икосаэдр  & 12  & 20   & 30  & 0         \\ \hline
    \end{tabular}
    \label{tab:b1}
\end{table}

\lipsum[1][1]

    \appendixsection{Приложение с формулой}

Формула~\ref{c1} также идёт с нумерацией приложения.

\begin{equation}
    2 + 2 = 4
    \label{c1}
\end{equation}

    \appendixsection{Приложение с листингом}

Рисунок~\ref{src:src3} с листингом в приложении.

\begin{figure}
\lstinputlisting[language=C++]{inc/example.cpp}
\caption{Пример использования листинга}
\label{src:src3}
\end{figure}
\end{document}
